\section{模型假设与符号说明}

\subsection{题设条件与估计假设}
张量大小、操作周期、数据依赖和存储位置以实际附件为准。计算图保持原语义，每个非 COPY 操作恰好执行一次；核心数、四条管道、缓存容量、同步延迟及带宽取题目配置。它们是约束模型的已知条件。

候选比较采用三项估计约定。第一，以当前图的计算工作量和依赖路径构造确定性的时间估计，不额外引入操作周期随机扰动。第二，问题一、二用任务或核心的汇总工作量近似资源拥堵，问题三将读写服务展开为事件；估计精度随模型层次增加，但不将任何一项代理量视为官方执行时间的恒等式。第三，核内驻留估计采用分析序列上的张量使用区间，实际空间复用及换出插入仍由固定执行机制确定。

这些约定使每次方案生成只依赖当前输入图、硬件配置和冻结参数。全局参数可以在开发阶段诊断、标定；正式生成方案时不读取历史答案或该图旧得分。已给定 100 图参与过开发，故全量成绩用于报告本题数据上的性能，独立输入测试另行列示。

\subsection{统一符号与计量单位}
原始图记为 $G=(T\cup O,E)$，非 COPY 操作集合为 $O_c$。本文用 $t$ 表示张量索引、$u,v$ 表示操作索引、$s$ 表示子图、$k$ 表示核心；物理时刻记为 $\theta$。张量大小统一用 byte，时间统一用 cycle，图中 MiB 按 $2^{20}$ byte 换算。

\begin{table}[htbp]
\centering
\caption{主要符号及其含义，局部辅助量在对应公式前定义}
\label{tab:notation}
\begin{tabular}{p{0.28\linewidth}p{0.59\linewidth}}
\hline
符号 & 定义与单位 \\
\hline
$s_t,c_u,p(u)$ & 张量字节数、操作周期、操作所属计算管道 \\
$G_c=(O_c,E_c)$ & 由原图生产--消费关系投影得到的操作依赖图 \\
$\mathcal C,C_j,K$ & 弱连通分量集合、第 $j$ 个分量及分量数 \\
$d_u,h,m_\ell$ & 操作拓扑深度、最大深度及第 $\ell$ 层操作数 \\
$N,k$ & 核心数与核心编号，$k=0,\ldots,N-1$ \\
$\mathcal S,\phi,a,\pi_k$ & 子图集合、操作归属、核心分配及核心 $k$ 的子图序列 \\
$P$ & 完整方案 $(\phi,a,\pi_0,\ldots,\pi_{N-1})$ \\
$W_j^q,L_k^q$ & 分量或核心在管道 $q\in\{M,V\}$ 上的工作量，cycle \\
$C_r,\beta,\beta_{\rm L2}$ & 存储容量 byte；DDR、L2 总带宽 byte/cycle \\
$T^{X}(P),D^{X}(P)$ & 场景 $X$ 中实际完工时间、官方额外搬运量 \\
$\widehat T,\widehat B,J$ & 时间估计、字节数估计及候选评分 \\
$T_{i,1},\overline S_N$ & 第 $i$ 图的整图单核参照、$N$ 核逐图加速比均值 \\
\hline
\end{tabular}
\end{table}

带帽变量来自求解模型，无帽的执行时间和搬运结果来自方案冻结后的验收。$\phi$ 与 $\pi_k$ 对应提交字段；$a$ 可由子图在核心序列中的位置恢复。该区分使建模解释、求解输出与结果统计保持一致。
